% Exemplo: título, fonte e nota na largura da tabela, com \IBGEtab (abntex2) ou threeparttable
% Compilar: latexmk -pdf tabela-fonte-nota.tex     Testado com: abntex2, threeparttable, TeX Live 2024+
% Cada trecho fica em uma página. As duas formas usam só a abntex2 e pacotes do CTAN; a classe
% abntex2-ufv faz o mesmo com o ambiente tabela (exemplos/abntex2-ufv/tabela-ufv.tex).
\documentclass[12pt, a4paper, brazilian]{abntex2}
\usepackage[T1]{fontenc}
\usepackage{lmodern}
%%BEGIN:tabelafontepreambulo
\usepackage{booktabs}                 % \toprule, \midrule, \bottomrule
\usepackage{threeparttable}           % título e notas na largura da tabela
%%END:tabelafontepreambulo

\begin{document}

A \autoref{tab:ibgetab} usa o \texttt{\textbackslash IBGEtab} da abntex2.

%%BEGIN:tabelaibgetab
\begin{table}[htbp]
\IBGEtab{\caption{Usuários ativos e tempo médio de sessão por plataforma}%
         \label{tab:ibgetab}}{%
  \begin{tabular}{lrr}
    \toprule
    \textbf{Plataforma} & \textbf{Usuários (un.)} & \textbf{Sessão (min)} \\
    \midrule
    Web     & 12.840 &  8,4 \\
    Android &  9.317 & 11,9 \\
    iOS     &  4.102 & 12,3 \\
    \bottomrule
  \end{tabular}%
}{\fonte{Elaborado pelo próprio autor.}%
  \nota{Dados fictícios, apenas para exemplo.}}
\end{table}
%%END:tabelaibgetab

\clearpage

A \autoref{tab:tpt} usa o ambiente \texttt{threeparttable}.

%%BEGIN:tabelatpt
\begin{table}[htbp]
  \centering
  \begin{threeparttable}
    \caption{Usuários ativos e tempo médio de sessão por plataforma}%
    \label{tab:tpt}
    \begin{tabular}{lrr}
      \toprule
      \textbf{Plataforma} & \textbf{Usuários (un.)} & \textbf{Sessão (min)} \\
      \midrule
      Web     & 12.840 &  8,4 \\
      Android &  9.317 & 11,9 \\
      iOS     &  4.102 & 12,3 \\
      \bottomrule
    \end{tabular}
    \fonte{Elaborado pelo próprio autor.}
    \nota{Dados fictícios, apenas para exemplo.}
  \end{threeparttable}
\end{table}
%%END:tabelatpt

\end{document}
